\chapter{OBJETIVOS}
\label{cap:objetivos}

\begin{borrador}
    \item \textbf{Fase:} 0 --- texto puente breve, redactable ya.
    \item Frase de enlace con el Capítulo~\ref{cap:introduccion}: la motivación y la brecha de usabilidad descritas allí se traducen aquí en objetivos concretos y verificables.
    \item Recordar en una frase el enfoque elegido (RAG con LLM local) sin repetir literalmente el resumen del Cap.~\ref{cap:introduccion}.
    \item Cerrar con una frase que anuncie la estructura de este capítulo: objetivo general, objetivos específicos, hipótesis de trabajo, restricciones y requisitos.
\end{borrador}

\section{Objetivo general}

Diseñar, implementar y evaluar un sistema conversacional capaz de responder, a
partir de información oficial y actualizada, preguntas de estudiantes
preuniversitarios sobre las titulaciones de la Escuela Politécnica Superior de
Jaén, empleando una arquitectura de recuperación aumentada por generación (RAG)
sobre un modelo de lenguaje de código abierto ejecutado en local.

\section{Objetivos específicos}

% Están redactados para poder responder sí/no a cada uno al terminar el TFG.
\begin{itemize}
    \item \textbf{OE1.} Construir una colección documental estructurada y
    trazable con los planes de estudio, las guías docentes y las salidas
    profesionales de la EPSJ, mediante un proceso de extracción web propio.
    \item \textbf{OE2.} Definir y evaluar experimentalmente una estrategia de
    fragmentación (\glr{chunking}{chunking}) de esa colección adecuada al dominio.
    \item \textbf{OE3.} Seleccionar, con criterios cuantitativos, el modelo de
    incrustaciones (\textit{embeddings}) y la base de datos vectorial del
    sistema.
    \item \textbf{OE4.} Implementar el \textit{pipeline} RAG completo sobre un
    modelo de lenguaje de código abierto desplegado en local.
    \item \textbf{OE5.} Evaluar el sistema con métricas objetivas de
    recuperación y de generación.
    \item \textbf{OE6.} Ofrecer el sistema a través de una aplicación web de
    chat accesible y sin registro.
\end{itemize}

\section{Hipótesis de trabajo}

% La hipótesis debe ser falsable: al evaluar el sistema se confirma o se refuta.
Es posible construir, con recursos de cómputo modestos y modelos de código
abierto, un sistema RAG que responda preguntas sobre la oferta académica de la
EPSJ con una fidelidad y una capacidad de recuperación suficientes para resultar
útil a un estudiante preuniversitario, sin incurrir en respuestas inventadas
sobre asignaturas o titulaciones de un modelo de lenguaje grande de internet.

\section{Restricciones del proyecto}

El proyecto lo desarrolla una sola persona, con la dedicación equivalente a los
créditos ECTS del Trabajo Fin de Grado y dentro del calendario de la
convocatoria correspondiente. El sistema debe funcionar con recursos de cómputo
domésticos, sin acceso a hardware especializado, lo que condiciona el tamaño de
los modelos que pueden emplearse.

Como limitación conocida del dominio (Sección~\ref{secc:alcance}), algunas asignaturas de cursos superiores
de las titulaciones de implantación reciente no tienen todavía guía docente
publicada.

\section{Requisitos}

\begin{borrador}
    \item \textbf{Fase:} 0 --- texto puente de una o dos frases antes de la primera subsección (norma del tutor: nunca pasar de un título a otro sin texto), anunciando que los requisitos se dividen en los del usuario y los del sistema.
\end{borrador}

\subsection{Requisitos de usuario}

\begin{borrador}
    \item \textbf{Fase:} 0 (IT-22) --- redactable ya.
    \item Perspectiva del estudiante preuniversitario, no del sistema. Candidatos: poder preguntar en lenguaje natural sin conocer la jerga universitaria (ECTS, menciones, tipos de asignatura); recibir información fiable de planes de estudio, guías docentes y salidas profesionales; poder comparar titulaciones entre sí; usar el sistema sin registro ni instalación; que las respuestas no inventen datos.
    \item Hacer cada requisito rastreable a un objetivo específico (OE1--OE6): esa trazabilidad es lo que el tribunal valora.
    \item Decidir el formato: tabla con identificador y prioridad (RU-01, RU-02\ldots con MoSCoW) o prosa. La tabla es más rastreable; recuerda la norma del tutor de formato uniforme en todas las tablas.
\end{borrador}

\subsection{Requisitos del sistema}

\begin{borrador}
    \item \textbf{Fase:} 1 (IT-59) --- \textbf{diferido}: NO se redacta en Fase 0. La especificación técnica (qué debe cumplir el software y el hardware) depende de decisiones de las Fases 1 y 2 que aún no se han tomado.
\end{borrador}